\section{Analysis for \cref{sec:optimizing_mechanisms}}

\subsection{Proof of~\cref{thm:dual_fn}}\label{app:ssec:dual-proof-multipliers-proof}
\begin{namedtheorem}[\cref{thm:dual_fn}]
Let a finite $\deltaset = \{\contrib_i\}_{i=1}^k$ be given, and assume that the vectors $\{\contrib_i\}_{i=1}^k$ span $\R^n$. Assume that $\bfA$ is full-rank, and for $\dualv \in \R^k$ define 
\[
\Hv = [\contrib_1, \dots, \contrib_k] \diag(\dualv)^{1/2}, \quad \bfU = \Hv\Hv\tp.
\]
Then, for Lagrange multipliers $\bfv$ such that the $\bfU$ is full-rank, the Lagrange dual function $g$ can be expressed in closed form in terms of the Lagrange multipliers: 
\begin{equation}
g(\dualv) := \inf_{\bfX\text{ is PD}} L(\bfX, \dualv) = 2\,\tr \left( \big(\bfU\hlf \AAT \bfU\hlf \big)\hlf \right) - \sum_{\contrib\in\deltaset} \dualv_\contrib
\end{equation}
\end{namedtheorem}
\begin{proof}

Since there is some finite set of vectors $\contrib \in \R^n$ specifying $\deltaset$, the supremum in~\cref{eq:constrained_problem} may be reduced to a maximum over these elements.%

Our problem then takes the form:
\begin{align}\label{eq:simplified_constrained_problem}
    \inf_{\bfX \text{ is PD}} \quad &\tr(\AAT \bfX\inv) \notag \\
    \text{s.t.} \quad & \contrib\tp \bfX \contrib \le 1, \ \forall \contrib \in \deltaset.
\end{align}

Recall that we have defined
$\Hv = [\contrib_1, \dots, \contrib_k] \diag(\dualv)\hlf$, and $\bfU = \Hv\Hv\tp$. Now, note:
 \begin{equation}\label{eq:Udef}
  \bfU = \sum_\contrib \dualv_\contrib \contrib \contrib\tp 
  = \bfH \diag(\dualv) \bfH\tp = \Hv \Hv\tp.
\end{equation}

Introducing Lagrange multipliers $\dualv_\contrib \ge 0$, for the problem \cref{eq:simplified_constrained_problem} we form the Lagrangian for positive-definite $\bfX$:
\begin{align}
  L(\bfX, \dualv) 
  &= \tr(\AAT\bfX^{-1})
     + \sum_\contrib \dualv_\contrib \left(\contrib\tp \bfX \contrib - 1\right) \\
  &= \tr(\AAT\bfX^{-1})
      + \tr \left( \Big(\sum_\contrib \dualv_\contrib \contrib \contrib\tp\Big) \bfX \right)
      - \sum_\contrib \dualv_\contrib \\
   &= \tr(\AAT\bfX^{-1})
      + \tr \left( \bfU \bfX \right) - \sum_\contrib \dualv_\contrib. \label{eq:genL}
\end{align}


For fixed $\dualv$, any finite minimizer of $L$ for positive-definite $\bfX$ must correspond to a zero of this Lagrangian's gradient. We then compute the gradient
\begin{equation}\label{eq:lagrangian_gradu}
  \frac{\partial L}{\partial \bfX} = -\bfX^{-1} \AAT \bfX^{-1} + \bfU.
\end{equation}
$\bfU$ and $\bfA$ are full-rank by assumption; therefore \cref{lem:Xsoln} is applicable, and \cref{eq:lagrangian_gradu} has a unique positive-definite zero (and indeed, the infimum in \cref{eq:simplified_constrained_problem} becomes a minimum):
\begin{equation}\label{eq:xoptu}
\bfX = \bfU\nhlf \big(\bfU\hlf \AAT \bfU\hlf \big)\hlf \bfU\nhlf.
\end{equation}
Note that~\cref{eq:lagrangian_gradu} also immediately implies that if $\bfU$ is not full-rank, then there is no finite positive-definite minimizer of $L$ in $\bfX$. Letting $g(\dualv) = \min_{\bfX} L(\bfX, \dualv)$ be the Lagrange dual function and plugging back into \cref{eq:genL}, we have 
\begin{align*}
g(\dualv) 
 &= \min_{\bfX \text{is PD}} \tr(\AAT\bfX^{-1})
      + \tr \left( \bfU \bfX \right) - \sum_\contrib \dualv_\contrib \\
 &= \min_{\bfX \text{is PD}} \tr(\bfX \bfU)
      + \tr \left( \bfU \bfX \right) - \sum_\contrib \dualv_\contrib 
      && \text{using \cref{eq:lagrangian_gradu}}\\
 &=  \min_{\bfX \text{is PD}} 2\,\tr \left( \bfU \bfX \right) - \sum_\contrib \dualv_\contrib \\      
 &=  2\,\tr \left( \bfU \bfU\nhlf \big(\bfU\hlf \AAT \bfU\hlf \big)\hlf \bfU\nhlf \right) - \sum_\contrib \dualv_\contrib
      && \text{using \cref{eq:xoptu}} \\
 &=  2\,\tr \left( \big(\bfU\hlf \AAT \bfU\hlf \big)\hlf \right) - \sum_\contrib \dualv_\contrib.
\end{align*}
\end{proof}

\subsubsection{Proof of~\cref{cor:gen_fixed_point}}
\begin{namedcorollary}[\cref{cor:gen_fixed_point}]
In the same setup as \cref{thm:dual_fn}, a maximizer of the dual $\vopt$ must satisfy:
\begin{align}
\vopt 
  &= \diagpart\left(\big(\Hvopt\tp \AAT \Hvopt\big)\hlf\right).
\end{align}
Moreover, the optimal value of the problem defined in \ref{eq:symmetrized_problem} is $\tr\left(\vopt\right)$.
\end{namedcorollary}
\begin{proof}
\
As noted in the remark after \cref{thm:dual_fn}, any optimal setting of the dual variables must be in the interior of a neighborhood in which the representation \cref{eq:dual_fn_rep} is valid. It is therefore permissible to differentiate this representation.


Differentiating, we find:

\begin{align*}
 \frac{\partial}{\partial \dualv_i} \tr \left( \big(\bfU\hlf \AAT \bfU\hlf \big)\hlf \right)
 &=  \frac{1}{2} \tr\left( \AAT \bfU\hlf (\bfU\hlf \AAT \bfU\hlf)\nhlf \bfU\nhlf \contrib_i \contrib_i\tp \right) \\
&= \frac{1}{2} \tr\left( \bfU\nhlf (\bfU\hlf \AAT \bfU\hlf) (\bfU\hlf \AAT \bfU\hlf)\nhlf \bfU\nhlf \contrib_i \contrib_i\tp \right) \\
&= \frac{1}{2} \tr\left( \bfU\nhlf (\bfU\hlf \AAT \bfU\hlf)\hlf \bfU\nhlf \contrib_i \contrib_i\tp \right) \\
&= \frac{1}{2} \tr\left(\contrib_i\tp \bfU\nhlf (\bfU\hlf \AAT \bfU\hlf)\hlf \bfU\nhlf \contrib_i \right)\\
&= \frac{1}{2}\contrib_i\tp \bfX \contrib_i,  \\
\end{align*}

by defining $\bfX = \bfU\nhlf \big(\bfU\hlf \AAT \bfU\hlf \big)\hlf \bfU\nhlf$ (recalling the usage of the symbol $\bfX$ in \cref{eq:xoptu}).

Therefore
\[
 \frac{\partial g(\dualv)}{\partial \dualv_i} =  \contrib_i\tp \bfX \contrib_i - 1,
\]

and we have the stated expression for the gradient of the dual function.

Now, at a maximizer of the dual function, this derivative must vanish.  An equivalent condition is $\diagpart(\bfH\tp \bfX \bfH) = \vec{1}$, and hence 
\begin{equation}
\tr(\bfU \bfX)
= \tr(\bfH \diag(\dualv) \bfH\tp \bfX ) 
= \tr(\diag(\dualv) \bfH\tp \bfX \bfH) = \sum_\contrib \dualv_\contrib,
\end{equation}
so at the optimum $\vopt$ in fact $g(\vopt) = \sum_\contrib \vopt_\contrib$, establishing the second claim of our result.


Again using the observation that $\diagpart(\Hv\tp \bfX \Hv) = \vec{1}$ and so 
$$\diagpart(\Hv\tp \bfX \Hv) = \diagpart(\diagv\bfH\tp \bfX \bfH) = \dualv.$$

Further, using the second claim of \cref{cor:reps_for_x}, we can take
\[
\bfX = \Hv\tpinv \big(\Hv\tp \AAT \Hv\big)\hlf \Hv\inv,
\]
and multiplying this by $\Hv\tp$ and $\Hv$ on the left and right respectively yields
\begin{align*}
\Hv\tp \bfX \Hv 
 &= \Hv\tp \Hv\tpinv \big(\Hv\tp \AAT \Hv\big)\hlf \Hv\inv \Hv 
 = \big(\Hv\tp \bfV \Hv\big)\hlf
\end{align*}
and so we conclude for the optimal Lagrange multiplier $\vopt$,
\begin{align}
\vopt 
  &= \diagpart\left(\big(\Hvopt\tp \AAT \Hvopt\big)\hlf\right).
\end{align}
\end{proof}

\subsection{Lemmas and Corollaries}\label{app:ssec:finite-minimizer}

\begin{lem}\label{lem:finite_minimizer}
The set of positive-definite $\bfX$ such that $\sup_{\contrib \in \deltaset} \contrib\tp\bfX\contrib \leq 1$ is bounded as a subset of $\R^\dimdim$ if and only if $\deltaset=\{\contrib\}$ spans $\R^n$.
\end{lem}

\begin{proof}
Suppose that $\deltaset$ spans $\R^n$. For a PSD matrix, a bound on the trace implies a bound on the elements; therefore it is sufficient to show that $\sup_{\contrib \in \deltaset} \contrib\tp\bfX\contrib \leq 1$ implies that the maximum eigenvalue of $\bfX$ is uniformly bounded for $\bfX$ PSD.

Take some set of vectors $\{\contrib_i\}_{i=1}^\dimension \in \deltaset$ which span $\R^\dimension$. Fix some representation

$$
\bfe_i = \sum_{j=1}^\dimension \alpha_{ij} \contrib_j,
$$

where $\bfe_i$ is the $i^{th}$ standard basis vector.

Take $\bfy$ of $\ell_2$ norm 1, and express:

$$\bfy = \sum_{i=1}^n \gamma_i \bfe_i.$$
Then for $\bfX$ satisfying our assumptions,

\begin{equation}
    \left|\bfy\tp\bfX\bfy\right| = \left|\sum_{i,j=1}^n \gamma_i\gamma_j \bfe_i\tp\bfX\bfe_j \right|= \leq n^2 \max_{1 \leq i, j \leq \dimension} |\gamma_i\gamma_j| \left|\bfe_i\tp\bfX\bfe_j \right|.
\end{equation}

Similarly, 

\begin{equation}
    \left|\bfe_i\tp\bfX\bfe_j \right| = \left|\sum_{k, l=1}^\dimension \alpha_{ik}\alpha_{jl} \contrib_k\tp\bfX\contrib_l\right| \leq  n^2 \max_{1 \leq k, l \leq \dimension} |\alpha_{ik}\alpha_{jl}| \left|\contrib_k\tp\bfX\contrib_l
    \right| \leq  n^2 \max_{1 \leq k, l \leq \dimension} |\alpha_{ik}\alpha_{jl}|,
\end{equation}

where the final inequality follows by the assumptions on $\bfX$.

Now, since the $\ell_2$ norm of $\bfy$ is 1, the orthogonality of the $\bfe_i$ imply that each $|\gamma_i\gamma_j|$ is at most 1. Therefore:

\begin{equation}
\left|\bfy\tp\bfX\bfy\right| \leq n^4 \max_{1 \leq i, j, k, l \leq \dimension} |\alpha_{ik}\alpha_{jl}|,
\end{equation}

and we have sufficiency of $\deltaset$ spanning $\R^n$.

For necessity, suppose $\deltaset$ does not span $\R^n$. Then there is some vector $\bfy \in \text{span}\left(\deltaset\right)^\perp$ of norm 1. Take any $\bfX$ such that $\sup_{\contrib \in \deltaset} \contrib\tp\bfX\contrib \leq 1$. Then, since $\bfy\tp\contrib = 0$ for all $\contrib \in \deltaset$, $\bfY := \bfX + \alpha \bfy$ satisfies the same set of inequalities for any $\alpha$.
\end{proof}

\begin{lem} \label{lem:Xsoln}
Let $\bfU, \bfV\in \Snpp$. Let $\bfU = \UL \UR$ be a factorization of $\bfU$ such that $\UR \bfV \UL$ is PSD, and the following equation defines a positive-definite matrix $\bfX$:
\begin{equation}\label{eq:Xsoln_genrep}
\bfX = \UR\pinv \big(\UR \bfV \UL \big)\hlf \UL\pinv.
\end{equation}
Then, this $\bfX$ solves the equation
\[
\bfX \bfU \bfX = \bfV
\quad \text{or equivalently} \quad
\bfU = \bfX\inv \bfV \bfX\inv.
\]
Moreover, this positive-definite solution $\bfX$ is unique.
\end{lem}

\begin{proof}

We will begin by showing that $\bfX$ as defined by \cref{eq:Xsoln_genrep} solves the equation $\bfX\bfU\bfX = \bfV$; then we will show that any two positive-definite representations of the form \cref{eq:Xsoln_genrep} are in fact identical.

Notice that the representation $\bfU = \UL \UR$ implies that $\rank(\UL) \geq \dimension$ and $\rank(\UR) \geq \dimension$. Therefore $\UL\UL\pinv = \bfI = \UR\pinv\UR$, as implied by the Moore definition of the Moore-Penrose pseudoinverse. So:

\begin{align*}
\bfX \bfU \bfX 
  & = \bfX \UL \UR \bfX \\
  &= \left(\UR\pinv \big(\UR \bfV \UL \big)\hlf \UL\pinv\right)
     \UL\UR
     \left(\UR\pinv \big(\UR \bfV \UL \big)\hlf \UL\pinv\right) \\
  &= \UR\pinv \big(\UR \bfV \UL \big)\hlf  \bfP_{R(\UL\pinv)}\bfP_{R(\UR)} \big(\UR \bfV \UL \big)\hlf \UL\pinv
  \end{align*}

We claim that $\big(\UR \bfV \UL \big)\hlf  \bfP_{R(\UL\pinv)} =  \bfP_{R(\UR)} \big(\UR \bfV \UL \big)\hlf = \big(\UR \bfV \UL\big)\hlf$. In both cases, this can be seen by multiplying on the left or right as appropriate by $\big(\UR \bfV \UL\big)\hlf$, and noting $\bfP_{R(\UR)} \UR = \UR$, $\UL \bfP_{R(\UL\pinv)} = \UL$. Since all the terms here are symmetric, the appropriate equality follows by uniqueness of the symmetric matrix square root. Therefore:

\begin{align*}
  \bfX \bfU \bfX  &= \UR\pinv \UR \bfV \UL \UL\pinv \\     
  &= \bfV.
\end{align*}

The uniqueness of a positive-definite $\bfX$ solving $\bfX\bfU\bfX = \bfV$ follows from the uniqueness of the usual matrix square root. Indeed, assume $\bfY$ positive-definite satisfies $\bfY\bfU\bfY = \bfV$. Then:

\[
\left(\bfU^{1/2} \bfY \bfU^{1/2}\right)^2 = \bfU^{1/2} \bfV \bfU^{1/2}
\]

Since the positive-definite square root is uniquely determined, $\bfU^{1/2} \bfY \bfU^{1/2}$ is uniquely determined. Since $\bfU$ is invertible, $\bfY$ is uniquely determined as well, and we have $\bfY = \bfX$.
\end{proof}



\begin{cor}\label{cor:reps_for_x}
Two particular instantiations of \cref{lem:Xsoln} are of interest. $\bfX$ as the matrix geometric mean of $\bfU\inv$ and $\bfV$ (taking $\UL = \UR = \sqrt{\bfU})$:
\begin{equation}\label{eq:Xsoln}
\bfX = \bfU\nhlf \big(\bfU\hlf \bfV \bfU\hlf \big)\hlf \bfU\nhlf,
\end{equation}
and assuming the representation $\bfU = \Hv \Hv\tp$:
\begin{equation}\label{eq:XsolnInH}
\bfX = \Hv\tppinv \big(\Hv\tp \bfV \Hv\big)\hlf \Hv\pinv.
\end{equation}
\end{cor}

\begin{proof}

 By positive-definiteness of $\bfU$ and $\bfV$, \cref{eq:Xsoln} is clearly positive definite; \cref{eq:XsolnInH} may be seen to be positive definite via the SVD of the pseudoinverses involved. Symmetry is again clear. Therefore both representations satisfy the assumptions of \cref{lem:Xsoln}.
\end{proof}

\newcommand{\hPi}{\widehat{\Pi}}
\newcommand{\posnegu}{\textbf{$\pm$-contrib}}
\newcommand{\Xnn}{\textbf{non-neg $\bfX \ge 0$}}
\newcommand{\XPinn}{\textbf{non-neg $\bfX_\hPi \ge 0$}}
\subsection{Impact of non-negativity constraints} \label{sec:nonneg}

We consider three variations of the constraints in  \cref{eq:symmetrized_problem}. In the first, we have the default constraints
\begin{equation}\label{eq:default-constraint}
\max_{\contrib \in \Dcorners^1_\Pi} \contrib\tp \bfX \contrib \le 1.
\end{equation}
This requires enumerating the $\assumedsep 2^\maxpart$ corners of $\Dcorners^1_\Pi$, which is tractable for the small problems we consider here.
Next, we consider the approach used for our experiments:
\begin{equation} \label{eq:nonneg-constraint}
\max_{\contrib \in \Dcorners^1_\Pi, \contrib \ge 0} \contrib\tp \bfX \contrib \le 1
\quad \text{and} \quad
\bfX \ge 0.
\end{equation}
Note here we only have $\assumedsep$ non-negative vectors $\contrib$, and so only $\assumedsep + \dimension^2$ constraints.

Finally, we consider and ``in between'' set of constraints. This is the minimal set of constraints that allows us to apply claim 2 of \cref{thm:multidim}. In order to define these constraints, let $\hPi = \big\{ (i ,j) | (i, j) \in \pi, i\ne j, \pi \in \Pi\big\}$, the set of pairs of distinct iterations in which the same user could possibly participate.  For Claim 2 of \cref{thm:multidim}, because we apply the theorem to sub-matrices $\bfC\idx{:}{\pi}$, a sufficient condition is $\bfX\idx{i}{j} \ge 0$ for all $(i, j) \in \hat{\Pi}$.  Then, the constraints become:
\begin{equation} \label{eq:pih-constraint}
\max_{\contrib \in \Dcorners^1_\Pi, \contrib \ge 0} \contrib\tp \bfX \contrib \le 1
\quad \text{and} \quad
\forall (i, j) \in \hPi, \ \bfX\idx{i}{j} \ge 0.
\end{equation}
Since $|\hPi| = \assumedsep \maxpart (\maxpart - 1)$, we have $\assumedsep +  \assumedsep \maxpart (\maxpart - 1) < \assumedsep + \dimension^2$ constraints.

In \cref{tab:negentries} we show results for a tiny problem:  $n=6$, $\maxpart=3$, and $\assumedsep=2$.  However, we have run additional experiments that align with the conclusions reported here for moderately larger problems (noting the \textbf{full} approach quickly becomes prohibitively expensive).
For the prefix-sum workload, the choice of additional constraints does not impact the total error, as in all cases $\bfX \ge 0$.  However, for the momentum workload (with momentum 0.95), we see small gaps, indicating that imposing the additional constraints does come at at a small (from 16.115 to 16.134) impact on the total error for this workload.

\begin{table}[H]
\begin{center}
\renewcommand{\arraystretch}{1.4}
\begin{tabular}{llS[table-format=2.3]S[table-format=2.3]S[table-format=2.3]}
\toprule
workload $\bfA$ & constraints & $\min_{i,j \in [\dimension]} \bfX\idx{i}{j}$ & $\min_{i,j \in \hPi } \bfX\idx{i}{j}$ &     {Error}  \\ [1.75ex]
\midrule
 prefix-sum & \cref{eq:default-constraint} & -0.000 & -0.000 &  6.461 \\
 prefix-sum &     \cref{eq:pih-constraint} & -0.000 & -0.000 &  6.461 \\
 prefix-sum &  \cref{eq:nonneg-constraint} & -0.000 & -0.000 &  6.461 \\
   momentum & \cref{eq:default-constraint} & -0.031 & -0.031 & 16.114 \\
   momentum &     \cref{eq:pih-constraint} & -0.015 & -0.000 & 16.131 \\
   momentum &  \cref{eq:nonneg-constraint} & -0.000 & -0.000 & 16.134 \\
\bottomrule
\end{tabular}
\caption{Comparison of matrix mechanisms for $n=6$, $\maxpart=3$, and $\assumedsep=2$; the momentum workload uses momentum $0.95$.  Error is the root-total-squared-error, the square root of $\calL(\bfB, \bfC)$ defined in \cref{eq:l_expression}.} \label{tab:negentries}
\end{center}
\end{table}
